% Part: sets-functions-relations
% Chapter: functions
% Section: kinds

\documentclass[../../../include/open-logic-section]{subfiles}

\begin{document}

\olfileid[pt]{sfr}{fun}{kin}
\olsection{Tipos de Funções}

\begin{explain}
Será útil introduzir uma espécie de taxonomia para alguns dos
tipos de funções que encontramos com mais frequência.

Para começar, podemos querer considerar funções com a propriedade de que
todo \emph{membro} do contradomínio é um valor da função. Tais
funções chamam-se !!{surjective}s, e podem ser esquematizadas como na
\olref{fig:surjective}.

\begin{figure}[H]
  \olasset{assets/diagrams/surjective.tikz}
  \caption{Uma função !!{surjective} tem todo !!{element} do
    contradomínio como valor.}
  \ollabel{fig:surjective}
\end{figure}
\end{explain}

\begin{defn}[Função !!^{surjective}]
Uma função $f \colon A \rightarrow B$ é \emph{!!{surjective}} se e só se $B$
é também a imagem de~$f$, isto é, para todo $y \in B$ existe pelo menos
um $x \in A$ tal que~$f(x) = y$, ou em símbolos:
\[
  (\forall y \in B)(\exists x \in A)f(x) = y.
\]
Chamamos a tal função !!a{surjection} de $A$ para $B$.
\end{defn}

\begin{explain}
Se quiser mostrar que $f$ é !!a{surjection}, tem de mostrar que
todo objeto no contradomínio de $f$ é o valor de $f(x)$ para alguma
entrada $x$.

Note-se que qualquer função \emph{induz} !!a{surjection}. Afinal,
dada uma função $f \colon A \to B$, seja $f' \colon A \to \ran{f}$
definida por $f'(x) = f(x)$. Como $\ran{f}$ é \emph{definida} como
$\Setabs{f(x) \in B}{x \in A}$, esta função $f'$ é garantidamente
!!a{surjection}
\end{explain}

\begin{explain}
Agora, qualquer função mapeia cada entrada possível para uma única saída. Mas
há também funções que nunca mapeiam entradas diferentes para a mesma
saída. Tais funções chamam-se !!{injective}s, e podem ser esquematizadas
como na \olref{fig:injective}.
\begin{figure}[H]
  \olasset{assets/diagrams/injective.tikz}
  \caption{Uma função !!{injective} nunca mapeia dois
    argumentos diferentes para o mesmo valor.}
  \ollabel{fig:injective}
\end{figure}
\end{explain}

\begin{defn}[Função !!^{injective}]
Uma função $f \colon A \rightarrow B$ é \emph{!!{injective}} se e só se, para
cada $y \in B$, há no máximo um $x \in A$ tal que~$f(x) = y$. Chamamos a
tal função !!a{injection} de $A$ para~$B$.
\end{defn}

\begin{explain}
Se quiser mostrar que $f$ é !!a{injection}, tem de mostrar que,
para quaisquer !!{element}s $x$ e $y$ do domínio de $f$, se $f(x)=f(y)$, então
$x=y$.
\end{explain}

\begin{ex}
A função constante $f\colon \Nat \to \Nat$ dada por $f(x) = 1$ não é
nem !!{injective}, nem !!{surjective}.

A função identidade $f\colon \Nat \to \Nat$ dada por $f(x) = x$ é
tanto !!{injective} quanto !!{surjective}.

A função sucessora $f \colon \Nat \to \Nat$ dada por $f(x) = x+1$
é !!{injective} mas não !!{surjective}.

A função $f \colon \Nat \to \Nat$ definida por:
\[
  f(x) =
  \begin{cases}
    \frac{x}{2} & \text{se $x$ é par} \\
    \frac{x+1}{2} & \text{se $x$ é ímpar.}
  \end{cases}
\]
é !!{surjective}, mas não !!{injective}.
\end{ex}

\begin{explain}
Muitas vezes, queremos considerar funções que são tanto
!!{injective}s como !!{surjective}s. Chamamos a tais funções
!!{bijective}s. Parecem-se com a função esquematizada na
\olref{fig:bijective}. !!^{bijection}s são também por vezes chamadas
\emph{correspondências biunívocas}, pois associam de modo único elementos
do contradomínio a elementos do domínio.
\begin{figure}[H]
  \olasset{assets/diagrams/bijective.tikz}
  \caption{Uma função !!{bijective} associa de modo único os elementos do
    contradomínio aos do domínio.}
  \ollabel{fig:bijective}
\end{figure}
\end{explain}

\begin{defn}[!!^{bijection}]
Uma função $f \colon A \to B$ é \emph{!!{bijective}} se e só se é tanto
!!{surjective} como !!{injective}. Chamamos a tal função
!!a{bijection} de $A$ para~$B$ (ou entre $A$ e~$B$).
\end{defn}

\end{document}
